\section{Chart-wise gluing on toric CY$_3$ (Stratum i)}
\label{sec:gluing:toric-cech}
\medskip

Stratum~i is the toric regime: every top-dimensional cone of the fan
contributes a $\C^3$-chart whose local critical Hall algebra is the
Schiffmann--Vasserot positive half $Y^+(\widehat{\fgl}_1)$. The gluing
object in this section is finite: after a DWR/Ran refinement of the
toric atlas and finite bounds $(N,r,L,m)$, the chart data form an
oriented Rees Hall coefficient system on the finite \v{C}ech/Ran nerve.
The output is the Maurer--Cartan element in the total convolution dg Lie
algebra; completion, vanishing-cycle realisation, Drinfeld doubling, and
vertex-algebra evaluation are separate operations. Thus a $\C^3$ chart
contributes $Y^+(\widehat{\fgl}_1)$, not $\cW_{1+\infty}$, and the
finite assembly below is not a compact CoHA. The simplest non-trivial
instance is the resolved conifold
\(
 \bY = \mathrm{Tot}(\cO_{\bP^1}(-1) \oplus \cO_{\bP^1}(-1))
\)
with its two-chart atlas $\bY = U_+ \cup U_-$. The multi-chart
generalisation indexes over a toric fan; the compact-$4$-cycle
obstruction separates zero-interior-point diagrams from diagrams
requiring McKay resolution; the zero-interior case isolates the finite
Rees layer available before realisation.

% ============================================================================
\subsection{\v{C}ech gluing of two charts: the conifold paradigm}
\label{sec:gluing-p2-conifold}
% ============================================================================

The obstruction resolved by a two-chart atlas is the absence of a global
$\C^3$-parametrisation: the resolved conifold $\bY$ is a rank-$2$ vector
bundle over $\bP^1$, not an affine variety. Any chart-wise
Schiffmann--Vasserot construction must name the transition datum that turns
two copies of $\C^3$ into $\bY$, then transport that datum through every
intermediate step; the $(3,0)$-form, the polyvector fields, the Hall
product, and the orientation line.

\subsubsection{The two-chart atlas}

Let $z \in \bP^1$ be the standard affine coordinate on the base, with
$\tilde z = z^{-1}$ on the complementary chart. On the total space
\(
 \bY = \mathrm{Tot}(\cO_{\bP^1}(-1) \oplus \cO_{\bP^1}(-1))
\)
introduce fibre coordinates $u_+, v_+$ in the trivialisation over the
$z$-chart and $u_-, v_-$ in the trivialisation over the $\tilde z$-chart.
The cover
\[
 U_+ \;=\; \{z \neq \infty\} \times \C^2 \;\cong\; \C^3,
 \qquad
 U_- \;=\; \{\tilde z \neq \infty\} \times \C^2 \;\cong\; \C^3,
\]
has overlap $U_+ \cap U_- = (\bP^1 \setminus \{0, \infty\}) \times \C^2
\cong \bC^\times \times \C^2$. A section of $\cO_{\bP^1}(-1)$ is a function
$f(z)$ on $U_+$ that transforms to $\tilde z f(\tilde z^{-1}) = z^{-1} f(z)$
on $U_-$. Writing $u_\pm, v_\pm$ as the fibre coordinates of the two $\cO(-1)$
summands one finds the transition
\begin{equation}\label{gl:eq:conifold-transition-chart}
 \tilde z \;=\; z^{-1},
 \qquad
 u_- \;=\; z\, u_+,
 \qquad
 v_- \;=\; z\, v_+,
\end{equation}
as the natural datum forced by the $\cO(-1) \oplus \cO(-1)$ twist. The
standard atlas of $\bP^1$ carries an orientation flip across
$z \leftrightarrow \tilde z$: the holomorphic $1$-form $dz = -\tilde z^{-2}\,
d\tilde z$ reverses sign through the change of chart, and this sign
propagates to the $(3,0)$-form computation of the next subsection.

\subsubsection{Jacobian of the transition on the $(3,0)$-form}

The Calabi--Yau $(3,0)$-form $\Omega_{\bY}$ is specified on $U_+$ by the
requirement that it trivialise the canonical bundle of $U_+ \cong \C^3$:
\[
 \Omega_{\bY}\big|_{U_+} \;=\; dz \wedge du_+ \wedge dv_+.
\]
To transport $\Omega_{\bY}$ to $U_-$ one computes the Jacobian of the
transition~\eqref{gl:eq:conifold-transition-chart} step by step. The exterior
derivatives are
\[
 dz \;=\; -\tilde z^{-2}\, d\tilde z,
 \qquad
 du_+ \;=\; d(z^{-1} u_-)
 \;=\; -z^{-2}\, u_-\, dz + z^{-1}\, du_-,
 \qquad
 dv_+ \;=\; -z^{-2}\, v_-\, dz + z^{-1}\, dv_-,
\]
so that
\[
 dz \wedge du_+ \;=\; dz \wedge z^{-1}\, du_-,
 \qquad
 dz \wedge du_+ \wedge dv_+ \;=\; dz \wedge z^{-1}\, du_- \wedge z^{-1}\, dv_-
 \;=\; z^{-2}\, dz \wedge du_- \wedge dv_-.
\]
Substituting $dz = -\tilde z^{-2}\, d\tilde z$ and the fibre Jacobian
$z^{-2} = \tilde z^{\,2}$:
\begin{equation}\label{eq:jacobian-calculation}
 \Omega_{\bY}\big|_{U_+}
 \;=\; dz \wedge du_+ \wedge dv_+
 \;=\; \underbrace{z^{-2}}_{\text{fibre twist}}\,
 \underbrace{dz}_{-\tilde z^{-2}\, d\tilde z}
 \wedge du_- \wedge dv_-
 \;=\; \underbrace{\tilde z^{\,2} \cdot (-\tilde z^{-2})}_{=\;-1}\,
 d\tilde z \wedge du_- \wedge dv_-
 \;=\; -\,d\tilde z \wedge du_- \wedge dv_-.
\end{equation}
The two factors $\tilde z^{\,2}$ (fibre twist, one $\tilde z$ per $\cO(-1)$
summand) and $-\tilde z^{-2}$ (cotangent coefficient from $dz$) cancel
exactly, leaving no residual $\tilde z^{\pm}$ and only the
base-orientation sign $-1$ of $\bP^1$. This is the statement that
$K_{\bY}$ is trivial; the transition cocycle on $U_+ \cap U_-$
collapses to a constant $-1$, with no $\tilde z$-dependence; and is
precisely the cancellation required by the Calabi--Yau condition
$c_1(\cO(-1) \oplus \cO(-1)) = -2H = -c_1(T_{\bP^1})$ (degrees of the
two $\cO(-1)$-summands sum to $-\deg K_{\bP^1} = -2$).

\subsubsection{The chart-wise Schiffmann--Vasserot CoHA}

Each chart $U_\pm \cong \C^3$ carries the Jordan-triple quiver with
potential
\[
 Q_{\mathrm{Jor}} : \;\; \underbrace{\bullet}_{0}
 \;\;\text{with three loops}\;\; x_\pm, y_\pm, z_\pm,
 \qquad
 W_\pm \;=\; \mathrm{tr}\bigl(x_\pm[y_\pm, z_\pm]\bigr).
\]
The critical CoHA of $(Q_{\mathrm{Jor}}, W_\pm)$ is the positive half of the
affine Yangian of $\widehat{\fgl}_1$ (Schiffmann--Vasserot 2012
\texttt{arXiv:1202.2756}; Tsymbaliuk 2014 \texttt{arXiv:1404.5240}),
\[
 \CoHA\bigl(U_\pm, W_\pm\bigr) \;\cong\; Y^+\bigl(\widehat{\fgl}_1\bigr),
\]
with generators $e_r^{(\pm)}, f_r^{(\pm)}, h_r^{(\pm)}$ indexed by $r \in
\Z_{\geq 0}$ and structure functions given by the RSYZ shuffle kernel
(Rapcak--Soibelman--Yang--Zhao 2020 \texttt{arXiv:1810.10402} \S5,
equation (5.3), specialised to the single-vertex Jordan-triple case).
At weight $(1,1)$ the shuffle kernel evaluates to
\begin{equation}\label{eq:structure-constant-1-1}
 \varphi_{11}(u_1; u_2) \;=\; \frac{1}{(u_1 - u_2 - \varepsilon_1)(u_1 - u_2 - \varepsilon_2)}
 \;\;\xrightarrow{\;u_2 \to u_1\;}\;\;
 \frac{1}{\varepsilon_1 \varepsilon_2},
\end{equation}
the structure constant that every downstream computation must reproduce.

\subsubsection{Fourier--Mukai transition kernel}

On the overlap $U_+ \cap U_- = \bC^\times \times \C^2$ the derived category
$D^b(\mathrm{Coh}(U_+ \cap U_-))$ carries the autoequivalence implementing
the transition~\eqref{gl:eq:conifold-transition-chart} at the level of sheaves.
Its integral kernel is
\begin{equation}\label{eq:FM-kernel-conifold}
 K_{+-} \;=\; \cO_\Delta \otimes \pi_1^*\,\cO_{\bP^1}(-1)
 \;\in\; D^b\bigl(\mathrm{Coh}\,(U_+ \cap U_-) \times (U_+ \cap U_-)\bigr),
\end{equation}
with $\Delta$ the diagonal and $\pi_1$ the projection onto the first factor.
The transform $T_{+-}(\cF) := \pi_{2,*}\bigl(\pi_1^*\cF \otimes K_{+-}\bigr)$
acts on polyvector fields
$\mathrm{PV}^\bullet(U_\pm) = \bigoplus_p \bigwedge^p T_{U_\pm}$ through the
chain rule applied to~\eqref{gl:eq:conifold-transition-chart} combined with the
$\cO_{\bP^1}(-1)$-twist.

\paragraph{Chain rule on the cotangent basis.}
Differentiate~\eqref{gl:eq:conifold-transition-chart} to obtain
\[
 du_+ \;=\; d(z^{-1} u_-) \;=\; -z^{-2}\, u_-\, dz + z^{-1}\, du_-,
 \qquad
 dv_+ \;=\; -z^{-2}\, v_-\, dz + z^{-1}\, dv_-.
\]
On the vertical (fibre-degree) sector; the sector entering the
Schiffmann--Vasserot Hall product; only the $z^{-1}\,du_-$ and
$z^{-1}\,dv_-$ terms survive (the $dz$-terms drop under pairing with the
Jordan-triple fibre polyvectors by~\eqref{eq:structure-constant-1-1}).
The cotangent multiplier is therefore $z^{-1}$ per fibre slot, giving
$z^{-(a+b)}$ on a fibre-form $du_+^a \wedge dv_+^b$.

\paragraph{Polyvector transport via $\Omega_\bY$-duality.}
The critical CoHA identifies polyvectors with differential forms via
Serre duality against the $(3,0)$-form $\Omega_\bY$:
\[
 \Lambda^p\, T_{\bY} \;\cong\; \Omega^{3-p}_\bY \otimes K_\bY^{-1}
 \;\cong\; \Omega^{3-p}_\bY
\]
using triviality of $K_\bY$ (equation~\eqref{eq:jacobian-calculation}).
Under this identification the fibre-polyvector $\partial_{u_+}^a
\partial_{v_+}^b$ corresponds (up to sign) to the fibre-form
$du_+^a \wedge dv_+^b$ contracted against the base $dz$-direction of
$\Omega_\bY$. Its transport under $T_{+-}$ is the cotangent transformation
$du_+^a \wedge dv_+^b = z^{-(a+b)}\, du_-^a \wedge dv_-^b$.

\paragraph{$\cO(-1)$-kernel twist.}
The factor $\pi_1^*\cO_{\bP^1}(-1)$ in $K_{+-}$ contributes a further
multiplier $z^{-1}$ tracking the base-$\bP^1$ line-bundle weight carried
by CoHA sections (the critical CoHA acts on $\varphi_W$-twisted sheaves
that inherit the $\cO(-1)$-weight of the superpotential
$W_\pm = \tr(x_\pm [y_\pm, z_\pm])$ along the compact base): a section
$g$ of $\cO_{\bP^1}(-1)$ expressed as $g(z)$ on $U_+$ reads
$\tilde z\, g(\tilde z^{-1}) = z^{-1}\, g(z)$ on $U_-$.

\paragraph{Net multiplier on polyvector sections.}
The CoHA-relevant sections are polyvector densities $f \cdot \partial^a
\partial^b$ with $f$ a section of $\cO(-1)$ (the superpotential line
bundle restricted to the overlap), not a scalar function. Combining the
cotangent chain rule (fibre-form multiplier $z^{-(a+b)}$) and the
$\cO(-1)$-kernel twist (scalar-section multiplier $z^{-1}$),
\[
 \underbrace{\bigl[z^{-(a+b)}\bigr]}_{\text{cotangent chain rule, fibre slots}}
 \;\cdot\;
 \underbrace{\bigl[z^{-1}\bigr]}_{\cO(-1)\text{-kernel twist, base weight}}
 \;=\;
 z^{-(1 + a + b)}.
\]
The Fourier--Mukai transport of a Jordan-triple fibre polyvector density
therefore reads
\begin{equation}\label{eq:FM-multiplier}
 T_{+-}^* :\;
 f(z, u_+, v_+)\, \partial_{u_+}^a\, \partial_{v_+}^b
 \;\longmapsto\;
 f(\tilde z^{-1}, \tilde z\, u_-, \tilde z\, v_-)\;
 z^{-(1 + a + b)}\;
 \partial_{u_-}^a\, \partial_{v_-}^b,
\end{equation}
with $f$ understood as a section of $\cO(-1)$ on the base. The
base-polyvector sector (powers of $\partial_z$) picks up the
base-orientation multiplier $(-\tilde z^{\,2})$ per $\partial_z$-slot; the
Jordan-triple selection rule~\eqref{eq:structure-constant-1-1} restricts
the Hall product to the fibre sector $\alpha = 0$ on which the base
orientation enters only through the global sign
of~\eqref{eq:jacobian-calculation}.

\subsubsection{Finite \v{C}ech/Ran Rees gluing}

The pair $(U_+, U_-)$ with overlap $U_+ \cap U_-$ is an ordinary
\v{C}ech spine, not a Weiss cover. Factorisation locality enters only
after a DWR/Ran refinement. Fix finite bounds
\[
 N\quad\text{(holomorphic jets)},\qquad
 r\quad\text{(renormalisation order)},\qquad
 L\quad\text{(charge)},\qquad
 m\quad\text{(Ran arity)}.
\]
Let $\mathfrak U_{\bY}$ be a DWR-good refinement of the two toric charts
by polydiscs subordinate to $U_+$, $U_-$, and $U_+\cap U_-$. A
$p$-simplex of the finite Ran nerve has the form
\[
 \sigma=(S,\{P_s\Subset U_{\epsilon_{s,0}}\cap\cdots\cap
 U_{\epsilon_{s,p}}\}_{s\in S}),
 \qquad
 \epsilon_{s,j}\in\{+,-\},
\]
with $|S|\leq m$ and charge bounded by $\Gamma_{\leq L}$. On a vertex
lying in $U_\pm$ the finite cyclic critical model is the Jordan triple
$(Q_{\mathrm{Jor}},W_\pm)$ truncated at $(N,r,L)$; on an edge lying in
$U_+\cap U_-$ it is transported by the Fourier--Mukai kernel
$K_{+-}$ and the multiplier~\eqref{eq:FM-multiplier}.

\begin{lemma}[Conifold \v{C}ech/Ran cyclic atlas]
\label{lem:conifold-cech-ran-cyclic-atlas}
For every finite bound $(N,r,L,m)$ the data above form a multiplicative
finite DWR/Ran cyclic critical atlas on $\mathfrak U_{\bY}$. Its
$0$-simplex Hall cohomology is the local positive half
\[
 H^0\!\left(\Hall^{\mathrm{Rees},\mathrm{or}}_{U_\pm}\!/\hbar\right)
 \;\cong\;
 Y^+(\widehat{\fgl}_1),
\]
and the edge restriction is the same positive half transported by
$K_{+-}$; no vertex of the atlas is a $\cW_{1+\infty}$-module.
\end{lemma}

\begin{proof}
The Jacobian calculation~\eqref{eq:jacobian-calculation} gives a
constant transition for the holomorphic volume form, hence transports
the cyclic pairing without a residual Laurent monomial. The chain rule
and the $\cO(-1)$ kernel give exactly the fibre multiplier
of~\eqref{eq:FM-multiplier}; face restriction of a polydisc in
$U_+\cap U_-$ is therefore the pullback of the vertex cyclic model by
the same Fourier--Mukai transition. On the triple faces of the refined
nerve the two possible restrictions are equal because the line-bundle
transition cocycle of $\cO(-1)$ is multiplicative and the diagonal
kernels compose by convolution.

For a disjoint Ran family
$\sigma_1\sqcup\sigma_2$ the finite potential is
$W_{\sigma_1}\boxplus W_{\sigma_2}$ and the cyclic form is the external
orthogonal sum. Thom--Sebastiani identifies the corresponding Rees Hall
complex with the completed tensor product of the two Rees Hall
complexes. The vertex computation is the Schiffmann--Vasserot
$\C^3$ theorem applied to the Jordan triple, so the local cohomology is
$Y^+(\widehat{\fgl}_1)$. The construction has used only critical Hall
polyvectors and Rees deformation; no vertex-algebra representation or
Drinfeld double has been supplied.
\end{proof}

Let
\[
 \mathfrak M_{\hCS,\Hall}^{N,r,L,m}(\mathfrak U_{\bY})
 =
 \Tot_{\sigma\in\mathsf N^{\Ran}_{\le m}(\mathfrak U_{\bY};
 \Gamma_{\le L})}
 \Hom_{\mathrm{cont}}\!\left(
 B\Obs_{\hCS}^{q,\le r,\le N}(\sigma)_{\le L},
 \Hall_{\sigma}^{\mathrm{Rees},\mathrm{or}}
 \right)
\]
with differential $D_{\mathrm{tot}}=d_{\Hall}\circ(-)
-(-1)^{|-|}(-)\circ d_{B\Obs}+\delta_{\Cech/\Ran}$ and convolution
bracket induced by the bar coproduct and the Rees Hall product.

\begin{proposition}[Conifold \v{C}ech/Ran Maurer--Cartan equation]
\label{prop:conifold-cech-ran-rees-mc}
The relative chart maps of
Lemma~\ref{lem:conifold-cech-ran-cyclic-atlas} assemble to
\[
 \Theta_{\bY}^{\mathrm{Rees},N,r,L,m}
 :=
 \sum_{0\leq p\leq m}
 \sum_{\sigma\in\mathsf N^{\Ran}_{p,\leq m}
 (\mathfrak U_{\bY};\Gamma_{\leq L})}
 \int_{\Delta^p}\theta_{\sigma}^{\mathrm{rel}}
 \in
 \mathfrak M_{\hCS,\Hall}^{N,r,L,m}(\mathfrak U_{\bY}),
\]
and this element satisfies
\[
 D_{\mathrm{tot}}\Theta_{\bY}^{\mathrm{Rees},N,r,L,m}
 +\frac{1}{2}
 [\Theta_{\bY}^{\mathrm{Rees},N,r,L,m},
  \Theta_{\bY}^{\mathrm{Rees},N,r,L,m}]
 =0.
\]
Thus the conifold atlas gives a finite oriented hCS-to-Rees-Hall
comparison on the \v{C}ech/Ran nerve. It does not identify the conifold
with $\cW_{1+\infty}$, with a compact CoHA, or with a Hall--Drinfeld
double.
\end{proposition}

\begin{proof}
On every simplex the cyclic homological perturbation transfer gives a
relative Hamiltonian $W_\sigma$ over $\Omega^\bullet(\Delta^p)$
satisfying the BV master equation. The relative Fourier--BV
identification rewrites this equation as
\[
 (d_{\Hall}+d_{\Delta^p})\theta_\sigma^{\mathrm{rel}}
 -\theta_\sigma^{\mathrm{rel}}d_{B\Obs}
 +\frac{1}{2}
 [\theta_\sigma^{\mathrm{rel}},\theta_\sigma^{\mathrm{rel}}]_{\mathrm{conv}}
 =0 .
\]
Integration over $\Delta^p$ converts the $d_{\Delta^p}$ term into the
alternating sum over faces by Stokes' formula. Face compatibility in
Lemma~\ref{lem:conifold-cech-ran-cyclic-atlas} identifies that boundary
sum with $\delta_{\Cech/\Ran}$ in $D_{\mathrm{tot}}$. Summing over the
finite nerve gives the displayed Maurer--Cartan equation. The last
non-identification follows because the finite Rees Hall complex precedes
vanishing-cycle realisation, Hopf pairing, completion, and vertex-algebra
evaluation.
\end{proof}

\paragraph{Orientation sign in the finite Rees product.}
The base-orientation flip in
$\Omega_{\bY}|_{U_-}=-d\tilde z\wedge du_-\wedge dv_-$ enters the
orientation line $o_\sigma$ of the edge Rees Hall complex. On cross-chart
Hall products the sign is the Koszul sign in the Thom--Sebastiani
transport, hence a coefficient of the Maurer--Cartan element above. A
super-Yangian presentation may be recovered only after the separate
realised quiver-Hall identification with the Klebanov--Witten conifold
quiver; it is not the output of the finite \v{C}ech/Ran equation.

\paragraph{Coefficient check for the $(1,1)$ term.}
The $(1,1)$ structure constant in the finite Rees Hall product equals
$(\varepsilon_1 \varepsilon_2)^{-1}$ through three independent
computations:
\begin{enumerate}[label=\emph{(\roman*)}, leftmargin=2em]
\item \emph{\v{C}ech two-chart gluing}: the limit $u_2 \to u_1$ of the
 chart-wise Jordan shuffle~\eqref{eq:structure-constant-1-1}.
\item \emph{Van den Bergh tilting}: the integral of the $\varphi_W$-twisted
 Euler class over $\mathrm{Ext}^1(S_0, S_1) = \C^2$, the rank-$2$
 Ext-space between the simple modules of the Klebanov--Witten quiver.
\item \emph{Klebanov--Witten super-shuffle}: the Jordan kernel residue at
 $z_1 = z_2$ of the direct super-shuffle algebra.
\end{enumerate}
Agreement across these computations pins $(\varepsilon_1 \varepsilon_2)^{-1}$
as the local coefficient carried by the finite Rees Hall comparison.

% ============================================================================
\subsection{Multi-chart toric fan generalisation}
\label{sec:gluing-p2-multi-chart}
% ============================================================================

The two-chart finite Rees assembly generalises to any local toric CY$_3$. The
obstruction resolved is the absence of a global affine chart: a toric
variety is covered by affine opens $U_\sigma$ indexed by the maximal
cones $\sigma$ of the fan, and the finite Rees Hall coefficient system is
assembled by running the \v{C}ech/Ran machinery along the face lattice of
the fan.

\subsubsection{Toric fan and chart decomposition}

Let $X = X_\Sigma$ be a smooth toric CY$_3$ with fan $\Sigma$ in the lattice
$N \cong \Z^3$. The Calabi--Yau condition is that the primitive generators
of all rays of $\Sigma$ lie in a common affine hyperplane, equivalently
that there exists a lattice vector $m \in M = N^\vee$ such that
$\langle m, v_\rho \rangle = 1$ for every ray generator $v_\rho$. A
natural choice places all ray generators at height $1$ along the
last coordinate; the toric diagram is then the $2$-dimensional polytope
obtained by projecting the rays to the plane $\{x_3 = 1\}$.

\paragraph{Chart structure at top cones.}
Each maximal (top-dimensional) cone $\sigma \in \Sigma(3)$ of a smooth
toric CY$_3$ is generated by exactly three rays
$\rho_1, \rho_2, \rho_3$ forming a $\Z$-basis of $N$, and determines an
affine chart
\[
 U_\sigma \;=\; \mathrm{Spec}\,\C[\sigma^\vee \cap M] \;\cong\; \C^3.
\]
Coordinates $(x_1^\sigma, x_2^\sigma, x_3^\sigma)$ on $U_\sigma$ are the
characters of the dual basis of $\sigma^\vee$. The Jordan-triple
potential on each chart is
\[
 W_\sigma \;=\; \mathrm{tr}\bigl(x_1^\sigma [x_2^\sigma, x_3^\sigma]\bigr),
\]
and the chart-wise CoHA is $\CoHA(U_\sigma, W_\sigma) \cong
Y^+(\widehat{\fgl}_1)$ by the Schiffmann--Vasserot theorem applied to each
$\C^3$ with the standard Jordan-triple.

\paragraph{Equivariant weights and the CY condition.}
The three-torus $T^3 = (\bC^\times)^3$ acts on each $U_\sigma \cong \C^3$
with equivariant weights $(\varepsilon_1^\sigma, \varepsilon_2^\sigma,
\varepsilon_3^\sigma)$. The Calabi--Yau constraint on $W_\sigma$ to be a
tri-vector of vanishing weight is
\[
 \varepsilon_1^\sigma + \varepsilon_2^\sigma + \varepsilon_3^\sigma \;=\; 0,
\]
which specialises the chart-wise weights to those of a local CY$_3$. Along
a compact direction (a ray $\rho$ in $\Sigma(1)$ whose dual edge in the
toric diagram is an internal edge) the equivariant parameter vanishes:
the two fixed points of the compact $\bP^1$-base on that edge would carry
opposite weights, obstructing global regularity of sections. This is the
mechanism that set $\varepsilon_3 = 0$ for the conifold base in
\S\ref{sec:gluing-p2-conifold}.

\begin{proposition}[Affine hCS BV chart and CY transition law]
\label{prop:toric-affine-hcs-bv-transition}
Let \(U_\sigma\simeq\C^3\) be a maximal-cone chart of a smooth toric
CY$_3\), with coordinates \(x_1^\sigma,x_2^\sigma,x_3^\sigma\), and
write
\[
 \Omega_X|_{U_\sigma}
 =
 \lambda_\sigma\,
 dx_1^\sigma\wedge dx_2^\sigma\wedge dx_3^\sigma,
 \qquad
 \lambda_\sigma\in\mathcal O(U_\sigma)^\times .
\]
For a metric Lie algebra \(\mathfrak g\), the affine-chart hCS BV
complex is
\[
 \cE_{\hCS,c}(U_\sigma,\mathfrak g)
 =
 \Omega_c^{0,\bullet}(U_\sigma,\mathfrak g)[1],
 \qquad
 Q_{\hCS}=\dbar ,
\]
with pairing and classical action
\[
 (\alpha,\beta)_\sigma
 =
 \int_{U_\sigma}\Omega_X\wedge\langle\alpha,\beta\rangle ,
 \qquad
 I_\sigma(\alpha)
 =
 \int_{U_\sigma}\Omega_X\wedge
 \left(
 \frac12\langle\alpha,\dbar\alpha\rangle
 +\frac16\langle\alpha,[\alpha,\alpha]\rangle
 \right).
\]
An element of Dolbeault degree \(q\) has BV cohomological degree
\(q-1\). The pairing is non-zero only on components with
\(q_1+q_2=3\), hence has degree \(-1\) on
\(\Omega_c^{0,\bullet}[1]\); the displayed superfield action means
taking the \((0,3)\)-component before integrating. On non-compact
charts the compact-support convention is understood after restriction
to the overlap on which the transition map is applied.
If \(U_\sigma\cap U_{\sigma'}\neq\varnothing\) and
\(\phi_{\sigma\sigma'}\) is the holomorphic transition map from
\(\sigma\)-coordinates to \(\sigma'\)-coordinates, then
\[
 \lambda_\sigma
 =
 (\lambda_{\sigma'}\circ\phi_{\sigma\sigma'})
 \det\!\left(
 \frac{\partial x^{\sigma'}_a}{\partial x^\sigma_b}
 \right)_{a,b=1}^3 .
\]
The transition map
\[
 T_{\sigma\sigma'}^{\BV}\colon
 \cE_{\hCS,c}(U_{\sigma'},\mathfrak g)|_{U_{\sigma}\cap U_{\sigma'}}
 \longrightarrow
 \cE_{\hCS,c}(U_{\sigma},\mathfrak g)|_{U_{\sigma}\cap U_{\sigma'}},
 \qquad
 \alpha\longmapsto\phi_{\sigma\sigma'}^*\alpha ,
\]
commutes with \(\dbar\), preserves the BV pairing and action, and
satisfies the \v{C}ech cocycle identity on triple overlaps. The same
Jacobian identity transports the finite Rees Hall divergence by
conjugacy:
\[
 T_{\sigma\sigma'}^{\PV}\Delta_{\omega_{\sigma'}}
 =
 \Delta_{\omega_\sigma}T_{\sigma\sigma'}^{\PV},
 \qquad
 T_{\sigma\sigma'}^{\PV}\iota_{dW_{\sigma'}}
 =
 \iota_{dW_\sigma}T_{\sigma\sigma'}^{\PV},
\]
whenever \(W_\sigma=\phi_{\sigma\sigma'}^*W_{\sigma'}\). Equivalently
\(T_{\sigma\sigma'}^{\PV}\) conjugates the finite Rees differential
\(\iota_{dW}+\hbar\Delta_{\omega}\). Here
\(\omega_\sigma=\Omega_X|_{U_\sigma}\) is the finite cyclic volume form.
\end{proposition}

\begin{proof}
The formula for \(\lambda_\sigma\) is the equality
\(\Omega_X|_{U_\sigma}=\phi_{\sigma\sigma'}^*(\Omega_X|_{U_{\sigma'}})\)
written in coordinates. Since \(\phi_{\sigma\sigma'}\) is holomorphic,
\(\phi_{\sigma\sigma'}^*\dbar=\dbar\phi_{\sigma\sigma'}^*\). The change
of variables formula and the displayed Jacobian identity give
\[
 \int_{U_{\sigma'}}\Omega_X\wedge\langle\alpha,\beta\rangle
 =
 \int_{U_\sigma}\Omega_X\wedge
 \langle\phi_{\sigma\sigma'}^*\alpha,
        \phi_{\sigma\sigma'}^*\beta\rangle ,
\]
and the same calculation applies to the quadratic and cubic terms of
\(I_\sigma\), using invariance of \(\langle-,-\rangle\). Thus the
degree-\((-1)\) BV pairing, bracket, and master equation are
transported without an extra density. On triple overlaps the coordinate
transitions compose, and the determinants multiply; hence the
transition maps form a \v{C}ech cocycle. The divergence operator
\(\Delta_{\omega}\) is characterised on polyvectors by
\(\mathfrak F_\omega(\Delta_\omega\xi)=d\,\mathfrak F_\omega(\xi)\),
equivalently by \(L_\xi\omega=(\Delta_\omega\xi)\omega\) for vector
fields. Since
\(\mathfrak F_{\omega_\sigma}T_{\sigma\sigma'}^{\PV}
=\phi_{\sigma\sigma'}^*\mathfrak F_{\omega_{\sigma'}}\), the same
Jacobian conjugates \(\Delta_\omega\); the equality
\(W_\sigma=\phi_{\sigma\sigma'}^*W_{\sigma'}\) gives the contraction
term. This is precisely the finite Rees compatibility used by
Lemma~\ref{lem:relative-fourier-bv-rees-identification}.
\end{proof}

\subsubsection{Gluing cocycle on the face lattice}

Two top cones $\sigma, \sigma'$ of $\Sigma(3)$ sharing a common face
$\tau = \sigma \cap \sigma' \in \Sigma(2)$ of codimension one define an
overlap $U_\sigma \cap U_{\sigma'} = U_\tau$. The face $\tau$ is a
two-dimensional cone, generated by two rays; the corresponding overlap is
\[
 U_\tau \;\cong\; \bC^\times \times \C^2,
\]
with the $\bC^\times$-factor encoding the direction transverse to $\tau$
within the full three-dimensional lattice. The transition between the chart coordinates
$(x_1^\sigma, x_2^\sigma, x_3^\sigma)$ and $(x_1^{\sigma'}, x_2^{\sigma'},
x_3^{\sigma'})$ follows the lattice change of basis between $\sigma^\vee$
and $(\sigma')^\vee$: when $\sigma'$ differs from $\sigma$ by a toric flop
sending ray $\rho_1$ to $\rho_1'$ across the shared face $\tau$ (generated
by $\rho_2, \rho_3$), the conifold-type transition is
\begin{equation}\label{eq:multi-chart-transition}
 x_1^{\sigma'} \;=\; (x_1^\sigma)^{-1},
 \qquad
 x_i^{\sigma'} \;=\; x_1^\sigma \cdot x_i^\sigma
 \;\; \text{for } i = 2, 3,
\end{equation}
in direct analogy with the conifold
transition~\eqref{gl:eq:conifold-transition-chart}. More generally, if the
toric flop twists the fibre $\rho_1$-direction with multiplicity $(k_2, k_3)$
along $\rho_2, \rho_3$ (equivalently, the fibre bundle over the shared
$\bP^1_\tau$ is $\cO(-k_2) \oplus \cO(-k_3)$), the transition reads
$x_i^{\sigma'} = (x_1^\sigma)^{k_i}\, x_i^\sigma$ on the $i$-th fibre
coordinate. The Fourier--Mukai kernel on the overlap is
\[
 K_{\sigma \sigma'} \;=\; \cO_\Delta \otimes \pi_1^* \cO_{\bP^1_{\tau}}(-1),
\]
the $\cO(-1)$-twist sourced by the compact $\bP^1_\tau$-base transverse
to $\tau$; the polyvector multiplier
$(x_1^\sigma)^{-(1 + k_2\, a + k_3\, b)}$ generalises~\eqref{eq:FM-multiplier}
with $a, b$ the fibre polyvector indices. The Jacobian
argument of \S\ref{sec:gluing-p2-conifold} carries over with $z$ replaced
by $x_1^\sigma$ on each codimension-one face, and with the fibre twist
exponents $(k_2, k_3)$ replacing the conifold's $(1, 1)$; the
Calabi--Yau condition $k_2 + k_3 = 2$ (sum of $\cO(-k_i)$-degrees equals
$-\deg K_{\bP^1_\tau} = -2$) enforces triviality of
$K_{X}|_{\bP^1_\tau}$.

\paragraph{The face lattice cocycle.}
The collection $\{K_{\sigma \sigma'}\}_{\sigma, \sigma' \text{ adjacent}}$
forms a cocycle on the face lattice of $\Sigma$: at any triple overlap
$U_\sigma \cap U_{\sigma'} \cap U_{\sigma''}$ (corresponding to a
codimension-two face $\tau = \sigma \cap \sigma' \cap \sigma''$), the
compatibility
\[
 K_{\sigma \sigma''} \;=\; K_{\sigma' \sigma''} \circ K_{\sigma \sigma'}
\]
holds up to canonical isomorphism, because each kernel is the pullback of
$\cO_{\bP^1}(-1)$ along a $\bP^1$-fibre direction and the fibre directions
compose additively in the toric fan. The cocycle class is the derivative
of the data on the fan's $1$-skeleton and is classified by the sheaf
cohomology group $H^1(\Sigma, \underline{\mathrm{Pic}}(\bP^1))$ on the fan
viewed as a CW complex.

\subsubsection{Galakhov--Li--Yamazaki bond factors}

Galakhov--Li--Yamazaki (2021 \texttt{arXiv:2106.01230}; Li--Yamazaki 2020
\texttt{arXiv:2003.08909} eq~(8.125)) present the multi-chart gluing data
via \emph{bond factors} $\varphi^{a \Rightarrow b}(u)$ indexed by ordered
pairs $(a, b)$ of vertices in the toric quiver. The convention is:
\begin{equation}\label{eq:gly-bond-factor}
 \varphi^{a \Rightarrow b}(u)
 \;=\;
 \frac{\prod_{I \in \mathrm{Arr}(a \to b)} (u + h^{a \to b}_I)}{%
 \prod_{J \in \mathrm{Rel}(a \to b)} (u + h^{a \to b}_J)},
\end{equation}
with $h^{a\to b}_I$ the equivariant weight of the $I$-th arrow from
vertex $a$ to vertex $b$, and $h^{a\to b}_J$ the equivariant weight of
the $J$-th Jacobi relation ($\partial W / \partial (\text{arrow})$) tying
arrows from $a$ into and out of a trip through $b$. The \emph{numerator
is arrow weights; the denominator is Jacobi-relation weights}; the
$(u + h)$-convention signs track the Drinfeld shifted-$R$-matrix
normalisation (Neguț 2015 \texttt{arXiv:1512.06473} eq~(1.6);
Tsymbaliuk 2014 \texttt{arXiv:1404.5240}).

\paragraph{Conifold bond factor from first principles.}
The conifold quiver (Klebanov--Witten 1998 \texttt{hep-th/9807080}) has
two vertices $\{0, 1\}$, two arrows $a_1, a_2 : 0 \to 1$ with equivariant
weights $h_1, h_2$, two arrows $b_1, b_2 : 1 \to 0$ with weights
$-h_1, -h_2$, and the Klebanov--Witten superpotential
$W_{\mathrm{con}} = \tr(a_1 b_1 a_2 b_2 - a_1 b_2 a_2 b_1)$ of
$T^2_{\mathrm{CY}}$-weight zero. For the bond $0 \Rightarrow 1$, the
numerator of~\eqref{eq:gly-bond-factor} enumerates the two $0 \to 1$
arrows with weights $h_1, h_2$, giving $(u + h_1)(u + h_2)$. The
denominator enumerates the two Jacobi relations tied to $0 \Rightarrow 1$,
namely the F-term ideal generators $\partial W / \partial b_1,
\partial W / \partial b_2$ read as elements of the weight-graded path
algebra at vertex $0$: the KW ideal admits a canonical generating pair
$(\mathfrak{r}_{\mathrm{sym}}, \mathfrak{r}_{\mathrm{asym}}) = \bigl(
a_1 b_1 + a_2 b_2, \; a_1 b_2 a_2 b_1 - a_1 b_1 a_2 b_2\bigr)$ of
$T^2_{\mathrm{CY}}$-weights $0$ (symmetric moment-map trace) and
$h_1 + h_2$ (Hessian-quadratic deformation direction) respectively,
matching the Negu\c{t}--Tsymbaliuk rank-$2$ shuffle presentation
(\texttt{arXiv:1404.5240} §4). Substituting into the GLY
formula~\eqref{eq:gly-bond-factor}:
\begin{equation}\label{eq:conifold-bond-factor}
 \varphi^{0 \Rightarrow 1}_{\mathrm{con}}(u)
 \;=\;
 \frac{(u + h_1)(u + h_2)}{u \cdot (u + h_1 + h_2)},
\end{equation}
matching Neguț's conifold shuffle kernel (Neguț 2015
\texttt{arXiv:1512.06473} eq~(1.6)) and Li--Yamazaki's
BPS-quiver-algebra master formula (\texttt{arXiv:2003.08909} eq~(8.125)).
At the Calabi--Yau torus-equivariant slice $\varepsilon_1 + \varepsilon_2
+ \varepsilon_3 = 0$, the conifold weights are $h_1 = \varepsilon_1$,
$h_2 = \varepsilon_2$, and $h_1 + h_2 = \varepsilon_1 + \varepsilon_2 =
-\varepsilon_3$; substituting into~\eqref{eq:conifold-bond-factor}:
\[
 \varphi^{0 \Rightarrow 1}_{\mathrm{con}}(u)
 \;=\;
 \frac{(u + \varepsilon_1)(u + \varepsilon_2)}{u\,(u + \varepsilon_1 + \varepsilon_2)}
 \;=\;
 \frac{(u + \varepsilon_1)(u + \varepsilon_2)}{u\,(u - \varepsilon_3)}.
\]
Both presentations are equivalent through the Calabi--Yau constraint
$\varepsilon_3 = -\varepsilon_1 - \varepsilon_2$.

\paragraph{Arrow-weight / relation-weight trace on the toric diagram.}
The arrow count $|\mathrm{Arr}(a \to b)|$ matches the number of toric
diagram edges joining the dual vertices of $a$ and $b$; the
Jacobi-relation count $|\mathrm{Rel}(a \to b)|$ matches the number of
toric triangles sharing the edge. For the conifold ($2$ arrows, $2$
relations per direction) the $2 + 2$ tally reflects the two-edge,
two-triangle diagram. For local $\bP^2$ the bond factors are labelled
by the three $\Z_3$-characters of the McKay quiver (Beilinson nine-arrow
structure, \S\ref{sec:gluing-p2-compact-cycle}); for the banana
threefold, by the three edges of the banana diagram; for the suspended
pinch point, by the pinch-point bond structure.

\subsubsection{Examples}

\paragraph{$\C^3/\Z_n$ (linear toric).}
The fan has $n$ maximal cones arrayed along the toric diagram edge, each
$\C^3$; the face lattice is a linear chain with $n - 1$ codimension-one
faces; the gluing cocycle is a chain of $n - 1$ FM kernels, each a
$\cO(-1)$-twist. The finite Rees atlas carries the $A_{n-1}$-type
McKay quiver-Hall coefficient system, matching the Beasley--Plesser
antisymmetric quiver on $\C^3/\Z_n$ after the orbifold realisation.

\paragraph{Banana threefold.}
Three ray clusters meeting at a common vertex produce a fan with three
maximal cones and three codimension-one faces. The finite \v{C}ech/Ran atlas
requires three FM kernels on a triangular face lattice; the finite Rees
coefficient system is a triangular gluing of three
Schiffmann--Vasserot $Y^+(\widehat{\fgl}_1)$ copies, governed by the
three-arrow banana quiver with cubic potential after quiver-Hall
realisation.

\paragraph{Suspended pinch point.}
The SPP singularity $\{xy = zw^2\} \subset \C^4$ has toric diagram the
trapezoid $\{(0,0), (2,0), (1,1), (0,1)\}$ with minimal triangulation
into three triangles sharing a common internal edge (the "pinch"). The
fan has three maximal cones and three codimension-one faces; the
\v{C}ech/Ran cocycle is defined on this triangulated trapezoid; the
finite Rees coefficient system carries the SPP quiver of Uranga (1999
\texttt{hep-th/9811004}) and Aganagic--Karch--L\"ust--Miemiec (1999
\texttt{hep-th/9903093}) after the realised quiver-Hall step.

% ============================================================================
\subsection{The compact $4$-cycle obstruction}
\label{sec:gluing-p2-compact-cycle}
% ============================================================================

The obstruction resolved here is subtle: a local toric CY$_3$ with a
compact divisor; a maximal cone whose dual vertex lies in the
\emph{interior} of the toric diagram; admits a chart-wise $\C^3$ cover,
but the boundary \v{C}ech/Ran Rees atlas of
\S\ref{sec:gluing-p2-multi-chart} misses the wheel term of the compact
divisor.

\subsubsection{Interior lattice points and compact divisors}

Let $X_\Sigma$ be a local toric CY$_3$ with toric diagram $\Delta_\Sigma$ a
lattice polygon in the plane $\{x_3 = 1\}$. The duality between toric
fans and toric diagrams identifies:
\[
 \text{rays of } \Sigma \;\longleftrightarrow\; \text{vertices of } \Delta_\Sigma,
 \qquad
 \text{maximal cones} \;\longleftrightarrow\; \text{internal triangulation triangles},
\]
and one distinguishes:
\begin{itemize}
\item \emph{boundary rays}: rays whose dual vertex lies on the boundary of
 $\Delta_\Sigma$; these correspond to non-compact divisors in $X$;
\item \emph{interior rays}: rays whose dual vertex lies in the interior
 of $\Delta_\Sigma$; these correspond to compact divisors, i.e.\ compact
 $4$-cycles in $X$.
\end{itemize}
The resolved conifold has no interior lattice point: its toric diagram is a
unit square with vertices $(0,0), (1,0), (0,1), (1,1)$ and no interior
lattice point, consistent with $\chi(\bY) = 2$ and no compact $4$-cycle.
Local $\bP^2 = \mathrm{Tot}(K_{\bP^2})$ has exactly one interior lattice
point: its toric diagram is the triangle with vertices $(1,0), (0,1),
(-1,-1)$ (the three non-compact divisor rays at height $1$) and the
interior lattice point $(0,0)$ at the origin; the ray $(0,0,1)$ attached
to this interior point generates the zero-section $\bP^2 \subset K_{\bP^2}$,
a compact $4$-cycle.

\subsubsection{Why chart-wise gluing misses the compact divisor}

The boundary \v{C}ech/Ran Rees atlas of \S\ref{sec:gluing-p2-multi-chart} has
three ingredients on each chart: the Schiffmann--Vasserot
$Y^+(\widehat{\fgl}_1)$, the equivariant torus action, and the FM kernel
on codimension-one overlaps. None of these detect the cohomology of a
compact divisor:

\begin{itemize}
\item The chart-wise polyvector complex $\mathrm{PV}^\bullet(U_\sigma)$ is
 concentrated in $(0, *)$-degree (holomorphic polyvector fields on
 $\C^3$), and has no contribution from the class of a compact surface.
\item The FM kernels $K_{\sigma \sigma'}$ transport polyvectors across
 codimension-one faces but do not generate new cohomology classes:
 they are sheaves of derived Morita data, not cycle classes.
\item The $T^3$-equivariant localisation concentrates all integrals at
 the fixed points of the torus action, i.e.\ the maximal cones; the
 contribution of a compact divisor is the \emph{sum over the fixed
 points contained in the divisor}, which must be computed by a separate
 procedure.
\end{itemize}
The Neguț wheel-condition corrections; the shuffle-algebra
relations forced by the presence of an interior lattice point; are
missed by the boundary \v{C}ech/Ran Rees atlas. For local $\bP^2$ with single
interior point at the origin of the triangulation, the correct shuffle
algebra is the $\Z_3$-McKay variant of $Y^+(\widehat{\fgl}_1)$ whose
kernel carries an additional wheel-condition multiplier tracking the
$\Z_3$-orbit of the interior lattice point (Neguț 2015
\texttt{arXiv:1512.06473}; Rapcak--Soibelman--Yang--Zhao 2020
\texttt{arXiv:2001.10549}). The direct multi-chart \v{C}ech/Ran Rees atlas of
\S\ref{sec:gluing-p2-multi-chart} picks up only the boundary contributions
to $\varphi_{a \Rightarrow b}$ and misses this wheel-condition
multiplier.

\subsubsection{Resolution via McKay/orbifold route}

The standard resolution of the obstruction is to pass to the orbifold
quotient: local $\bP^2 = [\C^3/\Z_3]^{\mathrm{res}}$ (the crepant
resolution of the McKay orbifold), and apply the Bridgeland--King--Reid
equivalence
\[
 D^b\bigl(\mathrm{Coh}(K_{\bP^2})\bigr) \;\simeq\; D^b\bigl(\mathrm{Coh}^{\Z_3}(\C^3)\bigr)
\]
(Bridgeland--King--Reid 2001 \texttt{arXiv:math/9908027} Thm.~1.2). The
right-hand category is a single $\C^3$-chart with $\Z_3$-equivariance;
the finite Hall model of this orbifold chart carries the Neguț wheel-condition
multipliers automatically, since the $\Z_3$-action on $\C^3$ produces
exactly the interior-point shifts at each orbit element. The compact
$4$-cycle cohomology is recovered as the $\Z_3$-fixed part of the
$\C^3/\Z_3$ inertia stack.

Alternatively, one may keep the multi-chart cover and enrich each
chart-wise shuffle kernel with a higher-rank Neguț correction: replace
$\varphi_{a \Rightarrow b}(u)$ by the shuffle kernel
$\varphi_{a \Rightarrow b}^{\mathrm{rank-}r}(u)$ of the rank-$r$ shuffle
algebra, with $r$ the number of interior lattice points of the toric
diagram. The two resolutions; McKay/orbifold route (\S\ref{sec:gluing:mckay}) and
higher-rank shuffle correction; are Morita-equivalent under the
Rapcak--Soibelman--Yang--Zhao toric theorem (2020
\texttt{arXiv:2001.10549} Thm.~B), but differ in their chart-wise
presentations.

\begin{proposition}[Compact-$4$-cycle obstruction]
\label{prop:compact-cycle-obstruction}
Let $X_\Sigma$ be a local toric CY$_3$ with toric diagram $\Delta_\Sigma$
carrying $r \geq 1$ interior lattice points. The direct multi-chart
\v{C}ech/Ran Rees atlas of \S\ref{sec:gluing-p2-multi-chart}, built only
from chart-wise single-vertex Jordan triples, presents the boundary
positive halves and misses the interior-lattice wheel terms. A finite
Rees atlas for the full toric quiver Hall layer must replace the local
chart data by either:
\begin{enumerate}[label=\emph{(\alph*)}, leftmargin=2em]
\item the $G$-equivariant Schiffmann--Vasserot data on $[\C^3 / G]$,
 where $G \subset \mathrm{SU}(3)$ is the McKay group whose crepant
 resolution $\widetilde{\C^3/G}$ recovers $X_\Sigma$ (McKay/orbifold
 route; applies when $\Delta_\Sigma$ is a McKay triangle, including
 $X_\Sigma = K_{\bP^2}$ with $G = \Z_3$);
\item the multi-vertex Neguț shuffle algebra $Y^+(Q_X, W_X)$ whose
 vertex set matches the full toric quiver of $X_\Sigma$ (interior and
 boundary lattice points together) and whose bond factors carry the
 wheel-condition multipliers tracking the interior-point positions
 (higher-rank shuffle route).
\end{enumerate}
The boundary multi-chart Rees atlas with only the chart-wise
$Y^+(\widehat{\fgl}_1)$ data therefore stops at the boundary Rees Hall
coefficient system.
\end{proposition}

\begin{proof}
The chart-wise single-vertex polyvectors generate only the
boundary-vertex portion of the full toric quiver shuffle algebra; the
interior-vertex portion is generated by compositions extending into the
$G$-McKay inertia (Bridgeland--King--Reid 2001 \texttt{arXiv:math/9908027}
Thm.~1.2). In the McKay route the finite critical chart is
$[\C^3/G]$, so the twisted sectors of the inertia stack supply exactly
the missing interior weights. In the higher-rank shuffle route the
finite critical chart is the full toric quiver $(Q_X,W_X)$; the
Neguț wheel factors are relations in the Rees Hall product and are
visible before any completion. Rapcak--Soibelman--Yang--Zhao (2020
\texttt{arXiv:2001.10549} Thm.~B) identifies these two presentations of
the same toric quiver shuffle data. The single-vertex atlas lacks both
the twisted sectors and the wheel relations, hence cannot produce the
interior-lattice terms.
\end{proof}

% ============================================================================
\subsection{Finite zero-interior Rees assembly}
\label{sec:gluing-p2-local-to-global}
% ============================================================================

The combination of \S\S\ref{sec:gluing-p2-conifold}--\ref{sec:gluing-p2-compact-cycle}
has a sharp zero-interior consequence. Boundary cones and codimension-one
faces supply all Rees Hall coefficients; no compact-divisor wheel term
remains to be added.

\begin{proposition}[Finite zero-interior toric Rees assembly]
\label{prop:finite-zero-interior-toric-rees-assembly}
Let $X_\Sigma$ be a smooth local toric Calabi--Yau threefold whose toric
diagram $\Delta_\Sigma$ has zero interior lattice points. Choose a
DWR-good refinement $\mathfrak U_\Sigma$ of its finite toric
$\C^3$-atlas and finite bounds $(N,r,L,m)$. Then:
\begin{enumerate}[label=\emph{(\arabic*)}, leftmargin=2em]
\item the chart Jordan triples $(Q_{\mathrm{Jor}},W_\sigma)$ on maximal
 cones, together with the Fourier--Mukai transitions on common faces,
 form a multiplicative finite DWR/Ran cyclic critical atlas on
 $\mathfrak U_\Sigma$;
\item the Stokes sum
\[
 \Theta_{X_\Sigma}^{\mathrm{Rees},N,r,L,m}
 =
 \sum_{0\leq p\leq m}
 \sum_{\sigma\in
 \mathsf N^{\Ran}_{p,\leq m}(\mathfrak U_\Sigma;\Gamma_{\leq L})}
 \int_{\Delta^p}\theta_\sigma^{\mathrm{rel}}
\]
 satisfies the \v{C}ech/Ran Maurer--Cartan equation in the finite
 convolution dg Lie algebra
 $\mathfrak M_{\hCS,\Hall}^{N,r,L,m}(\mathfrak U_\Sigma)$;
\item every $0$-simplex lying in a toric chart has local Hall cohomology
 $Y^+(\widehat{\fgl}_1)$, and the positive-half quiver presentation
 attached to the toric diagram is obtained only after the separate
 realised quiver-Hall identification.
\end{enumerate}
This constructs the finite oriented Rees Hall comparison on the
\v{C}ech/Ran nerve. It does not construct $\cW_{1+\infty}$, a compact
CoHA, a completed critical CoHA, or a Hall--Drinfeld double.
\end{proposition}

\begin{proof}
The fan has finitely many maximal cones and codimension-one faces. On
each maximal cone the local model is the Jordan triple on $\C^3$, whose
critical Hall cohomology is $Y^+(\widehat{\fgl}_1)$ by
Schiffmann--Vasserot. On a shared face the transition is the toric
lattice change of basis; its Fourier--Mukai kernel is the same
line-bundle transport as in~\eqref{eq:FM-kernel-conifold}, with
multiplier computed by~\eqref{eq:FM-multiplier} after replacing the
conifold exponents by the face exponents $(k_2,k_3)$. The Calabi--Yau
condition makes the holomorphic-volume Jacobian constant on the face,
so the cyclic pairing restricts compatibly.

On codimension-two intersections the toric line-bundle cocycle
condition makes the two composites of Fourier--Mukai kernels canonically
isomorphic. On disjoint Ran configurations the potentials add by
$W_{\sigma_1\sqcup\sigma_2}=W_{\sigma_1}\boxplus W_{\sigma_2}$, and
Thom--Sebastiani gives the Rees Hall product. These two compatibilities
are precisely the face and multiplicative axioms of a finite DWR/Ran
cyclic critical atlas.

The relative BV master equation on every simplex becomes, after the
relative Fourier--BV Rees identification,
\[
 (d_{\Hall}+d_{\Delta^p})\theta_\sigma^{\mathrm{rel}}
 -\theta_\sigma^{\mathrm{rel}}d_{B\Obs}
 +\frac12[\theta_\sigma^{\mathrm{rel}},\theta_\sigma^{\mathrm{rel}}]_{\mathrm{conv}}
 =0 .
\]
Integrating over $\Delta^p$ and summing over the finite nerve converts
the simplex boundary term to $\delta_{\Cech/\Ran}$, hence gives the
Maurer--Cartan equation for
$\Theta_{X_\Sigma}^{\mathrm{Rees},N,r,L,m}$. Zero interior lattice
points remove the Neguț wheel corrections described in
Proposition~\ref{prop:compact-cycle-obstruction}; the atlas therefore
contains exactly the boundary face coefficients and no compact-divisor
term. Rapcak--Soibelman--Yang--Zhao identify the realised toric quiver
shuffle algebra after the quiver-Hall realisation is supplied; that
identification is not part of the finite Rees construction proved here.
\end{proof}

The class of local toric CY$_3$ covered by
Proposition~\ref{prop:finite-zero-interior-toric-rees-assembly} (zero interior lattice
points in $\Delta_\Sigma$) includes:
\begin{itemize}
\item $\C^3$ (one chart, no overlaps, local positive half
 $Y^+(\widehat{\fgl}_1)$); toric diagram =
 triangle $\{(0,0), (1,0), (0,1)\}$, $0$ interior points;
\item the resolved conifold (two charts, one overlap,
 finite Rees comparison with Klebanov--Witten quiver-Hall realisation
 available after the separate realisation step); toric diagram = unit
 square $\{(0,0), (1,0), (0,1), (1,1)\}$, $0$ interior points;
\item the banana threefold (local model: three $\bP^1$'s meeting at two
 common nodes; three maximal cones, three overlaps, three-loop fan);
 local toric diagram with no interior lattice points;
\item the suspended pinch point $\mathrm{SPP} = \{xy = zw^2\} \subset
 \C^4$ (three maximal cones, three overlaps); toric diagram the
 trapezoid $\{(0,0), (2,0), (1,1), (0,1)\}$, $0$ interior lattice points;
\item generalised conifolds $X_{n, m} = \mathrm{Tot}(\cO(-n) \oplus
 \cO(-m) \to \bP^1)$ with $n + m = 2$ (two charts, one overlap with
 $\cO(-n) \oplus \cO(-m)$-twist); toric diagram an elongated
 quadrilateral with no interior lattice points.
\end{itemize}

The class outside the proposition comprises exactly those local
toric CY$_3$ with one or more interior lattice points. Among toric
smooth surfaces giving local CY$_3$, these are: local $\bP^2 = K_{\bP^2}$
(one interior point $(0,0)$), local $\bP^1 \times \bP^1 = K_{\bP^1
\times \bP^1}$ (one interior point), local Hirzebruch $\bF_n$ for
$n \geq 1$ (one interior point), and local toric del Pezzo
$dP_k = \mathrm{Bl}_k \bP^2$ for $k = 1, 2, 3$ (the only toric del
Pezzos beyond $\bP^2$; each with one interior point in its toric
diagram). Non-toric del Pezzos $dP_k$ for $k \geq 4$ fall outside the
toric framework entirely. These all carry compact $4$-cycles and must
pass through Proposition~\ref{prop:compact-cycle-obstruction}'s
McKay/orbifold or higher-rank shuffle resolutions. \S\ref{sec:gluing:mckay} takes up
this resolution, replacing the boundary chart-wise \v{C}ech/Ran atlas by the
BKR equivalence and the inertia-stack decomposition of the orbifold
route; in particular, the nine-arrow Beilinson quiver of local $\bP^2$
enters through the $[\C^3/\Z_3]$ chart and its $\Z_3$ skew-group algebra,
not through the boundary $\C^3$-chart atlas alone.

\bigskip

The finite Rees assembly of Stratum~i is determined by three data: the
face-lattice cocycle with values in
$H^1(\Sigma, \underline{\mathrm{Pic}})$; the bond-factor presentation of
Galakhov--Li--Yamazaki on the toric diagram edges; the zero-interior-point
criterion that separates the boundary-chart diagrams from those
requiring McKay resolution. \S\ref{sec:gluing:mckay} addresses the McKay route for
diagrams with one or more interior points, and \S\ref{sec:gluing:vdb} takes up derived
Morita gluing for the Van den Bergh tilting presentation common to
both. The master test case $K3 \times E$, which combines Stratum i with
Strata ii--iv, is resolved in \S\ref{sec:gluing:k3e} via the full four-stratum
decomposition.
